\chapter{Introduction}\label{ch:intro}

SCONTO-SVU est un modèle de simulation construit avec AnyLogic 8.9. Il représente une
chaîne logistique selon le modèle de référence SCOR~\cite{apics2017} et une lecture VSM
(\textit{Value Stream Mapping}). Chaque poste, acteur et commande y est un agent, et les
décisions de pilotage sont prises par des holons organisés en niveaux stratégique,
tactique et opérationnel.

Avant les travaux présentés ici, le modèle disposait de deux stratégies de production et
d'une politique de réapprovisionnement de type (s,Q). La stratégie MTS
(\textit{Make-To-Stock}) produit vers un stock de produits finis. La stratégie MTO
(\textit{Make-To-Order}) lance une production à partir d'une commande client. La politique
(s,Q) commande une quantité économique lorsque le stock projeté passe sous un point de
commande. Ces mécanismes fonctionnaient, mais trois problèmes restaient ouverts.

Le premier concerne les engagements matière. Une matière physiquement présente pouvait être
lue comme disponible par deux commandes à la fois. Le deuxième concerne la visibilité : une
matière commandée mais non encore reçue n'était pas représentée comme une quantité engagée,
et elle n'apparaissait qu'au moment de la réception. Le troisième concerne la coordination :
une commande MTO pouvait déclencher un approvisionnement alors que le stock en transit le
couvrait déjà.

Le travail décrit ici répond à ces trois problèmes. Il introduit une réservation explicite
de matière pour chaque commande, puis un noyau DDMRP (\textit{Demand Driven Material
Requirements Planning}) de gestion matière, avec une représentation unique des
approvisionnements en transit. Il règle enfin la coexistence entre ce noyau et la politique
(s,Q) historique.

La portée doit être énoncée clairement. Ce travail implémente un \emph{noyau DDMRP} de
gestion matière. Il ne reproduit pas la méthodologie DDMRP complète : le positionnement
stratégique des stocks (\textit{Strategic Inventory Positioning}) n'est pas traité. Les
paramètres des buffers sont expérimentaux. Les jeux utilisés pour la validation sont des
jeux de test, et ne sont pas des paramètres réels de ZENER.

Le document suit une progression volontaire, de la distinction MTS / MTO jusqu'à la
validation expérimentale. La phrase qui résume l'architecture est la suivante :
\begin{quote}
\itshape DDMRP anticipe, Open Supply mémorise, MTO alloue, la réservation sécurise,
Make consomme et Deliver clôture.
\end{quote}
Chaque verbe correspond à une partie du simulateur (\cref{fig:chaine}).

\begin{figure}[H]
  \centering
  \begin{tikzpicture}[font=\small, >=Stealth,
    stp/.style={draw=ink, rounded corners=3pt, minimum width=2.1cm, minimum height=1.25cm,
                 align=center, fill=white}]
    \node[stp, fill=ddmrpgrn!25] (a) {DDMRP\\anticipe};
    \node[stp, fill=osc!22, right=0.55cm of a] (b) {Open Supply\\mémorise};
    \node[stp, fill=mtoc!25, right=0.55cm of b] (c) {MTO\\alloue};
    \node[stp, fill=resc!25, right=0.55cm of c] (d) {Réservation\\sécurise};
    \node[stp, fill=mtsc!20, right=0.55cm of d] (e) {Make\\consomme};
    \node[stp, fill=ink!10, right=0.55cm of e] (f) {Deliver\\clôture};
    \foreach \x/\y in {a/b,b/c,c/d,d/e,e/f} \draw[->, thick] (\x) -- (\y);
  \end{tikzpicture}
  \caption{Schéma central : DDMRP anticipe, Open Supply mémorise, MTO alloue, la réservation
    sécurise, Make consomme et Deliver clôture.}
  \label{fig:chaine}
\end{figure}

